\documentclass[12pt]{article}
\usepackage{amsmath}
\usepackage{amssymb}
\usepackage{geometry}
\geometry{margin=1.4in}
\usepackage{hyperref}
\usepackage{setspace}
\setstretch{1.4}
\usepackage{amsthm}
\newtheorem{conjecture}{Conjecture}

\newcommand{\Sha}{\mathrm{Sh}}

\title{One Wall}
\author{Diego Rinc\'on\\
\small Independent}
\date{}

\begin{document}

\maketitle

\begin{abstract}
The five extension walls are not five different problems.
They are five encounters with the same problem:
the local-to-global passage.
This is an observation, not a proof.
\end{abstract}

\section{The Pattern}

Five of the six Millennium Problems involve a wall.
The wall has been named in each case separately,
in the language of each field:
Hadamard product control for RH,
continuum limit for Yang--Mills,
enstrophy closure for Navier--Stokes,
complex regularity for Hodge,
Euler system rank for BSD.

These names are correct.
But they are local names.
The same wall has a global name:
the local-to-global passage.

Every extension wall is a failure to pass
from local data to global structure.

\section{The Five Encounters}

\subsection*{Riemann Hypothesis}

The zeta function factors locally:
$$\zeta(s) = \prod_p \frac{1}{1 - p^{-s}}.$$
Each factor is a local Euler factor at a prime.
The zeros of $\zeta(s)$ are global objects:
they encode the distribution of primes via the explicit formula
$$\psi(x) = x - \sum_\rho \frac{x^\rho}{\rho} - \log(2\pi) - \tfrac{1}{2}\log(1 - x^{-2}).$$
The Riemann Hypothesis is a statement about global zero structure.
The local Euler product does not determine the location of the zeros.
The Beurling--Nyman criterion recasts RH as a global $L^2(0,1)$ statement.
The Connes approach works in the global ad\`ele class space $\mathbf{A}_\mathbf{Q}/\mathbf{Q}^*$.
Both approaches are manifestly global.

\textit{The wall:} there is no known path from the local Euler factors
to the global zero distribution.

\subsection*{Yang--Mills}

The Yang--Mills action is local:
$$S[A] = \frac{1}{4g^2}\int_{\mathbf{R}^4} \mathrm{tr}(F_{\mu\nu}F^{\mu\nu})\,d^4x.$$
It integrates local curvature.
The lattice construction builds a local theory: on each plaquette,
a compact group element.
The mass gap $\Delta > 0$ is a global spectral statement:
the Hamiltonian $H$ has no spectrum in $(0,\Delta)$.
The continuum limit --- lattice spacing $a \to 0$ ---
is the passage from the local lattice structure
to the global quantum field theory.
The gap must survive this passage.

\textit{The wall:} the local lattice theory is well-defined;
the global continuum limit is not proved.

\subsection*{Navier--Stokes}

The Navier--Stokes equations are a local PDE:
$$\partial_t u + (u\cdot\nabla)u = \nu\Delta u - \nabla p, \quad \nabla\cdot u = 0.$$
Vortex stretching is a local mechanism:
$\frac{D\omega}{Dt} = (\omega\cdot\nabla)u + \nu\Delta\omega$.
In two dimensions, the global conservation of enstrophy
$X(t) = \tfrac{1}{2}\int |\omega|^2\,dx$
closes the local-to-global passage:
$\frac{d}{dt}X \leq CX$ in 2D, giving global regularity by Gr\"onwall.
In three dimensions, there is no enstrophy conservation;
the local vortex-stretching term drives
$\frac{d}{dt}X \leq \frac{C}{\nu^3}X^3$,
a cubic inequality that cannot be closed.

\textit{The wall:} local vortex stretching in 3D
cannot be controlled globally.
The 2D proof works because 2D has an exact local-to-global bridge;
3D does not.

\subsection*{Hodge}

Differential forms are local objects.
Cohomology is global:
$H^{p,p}(X, \mathbf{Q})$ depends on the global geometry of $X$.
The Hodge theorem is a local-to-global theorem:
every de~Rham cohomology class has a unique harmonic representative.
The Hodge conjecture asserts more:
every rational class of type $(p,p)$ is a linear combination of
classes of algebraic cycles $[Z]$.
The Complex Regularity Conjecture (the wall) states:
a globally calibrated current $T$ of type $(p,p)$
with $[T] \in H^{p,p}(X,\mathbf{Q})$
need not be locally complex-structured.

\textit{The wall:} a global cohomological condition
does not imply local algebraic structure.
The local-to-global passage runs the wrong way:
the global class is known; the local structure is what must be derived.

\subsection*{Birch--Swinnerton-Dyer}

The $L$-function $L(E,s) = \prod_p L_p(E,s)$
is built from local Euler factors.
The rank $r = \mathrm{rank}\,E(\mathbf{Q})$ is a global invariant.
BSD asserts $\mathrm{ord}_{s=1} L(E,s) = r$.
The Euler system method (Kolyvagin) constructs global arithmetic classes
from local cohomology data.
For $r \leq 1$, Heegner points provide the required class.
For $r \geq 3$, the GGZ$_r$ conjecture requires $r$ independent classes;
only one construction is known.

\textit{The wall:} the local Euler factors encode the global rank by hypothesis;
deriving $r \geq 3$ independent global classes from local data is open.

\section{The Sixth Problem}

P versus NP does not fit this pattern.

It is not a failure to pass from local to global.
It is a failure to know what kind of argument to make.
The three barriers (relativization, natural proofs, algebrization)
have eliminated the known methods.
What remains is not a wall at the boundary of a local-to-global passage
but a void where the method must live.

P versus NP is a barrier wall, not an extension wall.
The local-to-global observation applies to the other five.

\section{What the Pattern Means}

The local-to-global problem is the central problem of modern mathematics.

Class field theory: local class field theory at each prime
$\to$ global reciprocity (Artin map).
Hasse--Minkowski: local solutions at $\mathbf{R}$ and each $\mathbf{Q}_p$
$\to$ global rational solution (for quadratic forms).
The Langlands program: local Langlands at each prime
$\to$ global automorphic representation.
Bloch--Kato: local conditions on Galois representations
$\to$ global $L$-function behavior.

The arithmetic geometry cluster ---
Langlands, Shimura, perfectoid, anabelian,
deformation theory, Iwasawa, Bloch--Kato, motives ---
is the accumulated machinery for local-to-global problems.

BSD and RH are both about $L$-functions.
A global $L$-function is a product of local factors.
The question is always: what does the product know
that the factors do not?

This is the question the body has been circling.
The five walls are five faces of it.

\section{A Conjecture}

\begin{conjecture}[One Wall]
There exists a unified local-to-global principle
from which the five extension walls
are special cases.
A proof of any one of the five
that does not use such a principle
will fail to generalize.
A proof that does use such a principle
will, with modification,
yield progress on the others.
\end{conjecture}

This is speculative.
It is not known.
It may be wrong.

But the pattern is real:
five problems, five walls,
all at the same passage.
The local data is under control.
The global structure is not.

If the pattern is not coincidence ---
if the five walls are one wall ---
then the map of the walls
is a map of one obstacle,
seen from five positions.

The body is five maps of one mountain.

\end{document}
